\chapter{Multi-Result Supercompilation (MRSC)}
\label{chap:mrsc}

\section{The Heuristic Dilemma in Metacomputation}
\label{sec:mrsc:dilemma}

Traditional supercompilers make hard heuristic choices at each transformation step:
\begin{itemize}
    \item Should an expression be inlined or driven as an uninterpreted function call?
    \item Should a whistle trigger generalization or term decomposition (splitting)?
    \item Which ancestor configuration in the path should be chosen for folding?
\end{itemize}
A suboptimal choice early in process tree exploration can lead to severe residual code bloat, miss crucial deforestation opportunities, or trap the compiler in complex knot topologies.

\section{Mitchell--Klyuchnikov Configuration Hypergraphs}
\label{sec:mrsc:hypergraphs}

To eliminate heuristic fragility, Mitchell and Klyuchnikov formulated \emph{Multi-Result Supercompilation (MRSC)} \citep{klyuchnikov2012practical}. Instead of constructing a single process tree, MRSC constructs a \emph{configuration hypergraph} ($\mathcal{H} = \langle V, E \rangle$) representing all legally permissible supercompilation derivations.

Implemented in \texttt{src/mir/supercompiler/mrsc.rs}:
\begin{lstlisting}[language=Rust, caption={The MRSC Hypergraph structures in NumLang.}]
pub struct MrscHypergraph {
    pub nodes: Vec<MrscNode>,
    pub hyperedges: Vec<MrscHyperedge>,
}

pub struct MrscHyperedge {
    pub source: MrscNodeId,
    pub targets: Vec<MrscNodeId>,
    pub label: TransformationRule,
}
\end{lstlisting}

\section{The Four-Dimensional MRSC Cost Model}
\label{sec:mrsc:cost_model}

Once the hypergraph is populated, NumLang evaluates candidate residual programs against an explicit multi-dimensional cost function:
\begin{equation}
\mathcal{C}(P) = w_{\text{step}} \cdot S(P) + w_{\text{alloc}} \cdot A(P) + w_{\text{code}} \cdot C(P) + w_{\text{reg}} \cdot R(P)
\end{equation}
where:
\begin{itemize}
    \item $S(P)$ is the estimated dynamic step count along hot paths.
    \item $A(P)$ is the static number of heap allocations.
    \item $C(P)$ is the residual machine instruction count (binary footprint).
    \item $R(P)$ is the estimated register pressure (live variable count).
\end{itemize}

The extraction procedure \texttt{extract\_pareto\_optimal} traverses the hypergraph via dynamic programming, extracting the residual that minimizes $\mathcal{C}(P)$.

\section{The Phase 53 Exhaustive MRSC Oracle}
\label{sec:mrsc:oracle}

When maximum optimization is required, the CLI flag \texttt{--mrsc-exhaustive} activates the Phase 53 \emph{Exhaustive MRSC Oracle} (\texttt{src/mir/supercompiler/mrsc\_oracle.rs}). 

The oracle replaces bounded graph pruning with an unbounded Iterative Deepening Depth-First Search (IDDFS). By exhaustively exploring alternative generalization frontiers, the oracle routinely discovers residual configurations that are strictly smaller and faster than online greedy supercompilation. Optimal residuals discovered by the oracle are automatically committed to the L2 persistent cache for future compilations.

\section{The Exhaustive IDDFS Oracle Implementation (Phase 53)}
\label{sec:mrsc:oracle_code}

Listing~\ref{lst:mrsc_oracle_code} presents the implementation of the Iterative Deepening Depth-First Search oracle from \texttt{src/mir/supercompiler/mrsc\_oracle.rs}:

\begin{lstlisting}[language=Rust, caption={Exhaustive MRSC Oracle Engine in NumLang (\texttt{src/mir/supercompiler/mrsc\_oracle.rs}).}, label={lst:mrsc_oracle_code}]
use crate::mir::lower::MirFunction;

/// Configuration for the MRSC Iterative Deepening Depth-First Search (IDDFS) Oracle.
#[derive(Debug, Clone)]
pub struct IddfsOracleConfig {
    /// Minimum search depth bound (default: 2).
    pub min_depth: usize,
    /// Maximum search depth bound (depth >= 20 for exhaustive exploration, default: 24).
    pub max_depth: usize,
    /// Depth increment step per deepening iteration (default: 2).
    pub step_depth: usize,
    /// Objective function for selecting from the Pareto frontier.
    pub objective: MrscObjective,
    /// Whether exhaustive search mode is enabled.
    pub exhaustive: bool,
}

impl Default for IddfsOracleConfig {
    fn default() -> Self {
        Self {
            min_depth: 2,
            max_depth: 24,
            step_depth: 2,
            objective: MrscObjective::Balanced,
            exhaustive: true,
        }
    }
}

/// A candidate residual program discovered during IDDFS exploration.
#[derive(Debug, Clone)]
pub struct OracleCandidate {
    pub depth: usize,
    pub residual: MirFunction,
    pub tree: ProcessTree,
    pub cost: MrscCostVector,
    pub fitness_score: f64,
}

/// Pareto frontier over the 4-dimensional cost model.
#[derive(Debug, Clone, Default)]
pub struct OracleParetoFrontier {
    pub candidates: Vec<OracleCandidate>,
    pub max_depth_evaluated: usize,
    pub total_evaluated: usize,
}

impl OracleParetoFrontier {
    pub fn new() -> Self {
        Self {
            candidates: Vec::new(),
            max_depth_evaluated: 0,
            total_evaluated: 0,
        }
    }

    /// Inserts a candidate into the 4D Pareto frontier.
    /// Rejects the candidate if dominated by any existing member.
    /// Removes any existing members dominated by the new candidate.
    pub fn insert(&mut self, candidate: OracleCandidate) -> bool {
        if self.candidates.iter().any(|c| c.cost.dominates(&candidate.cost)) {
            return false;
        }
        self.candidates.retain(|c| !candidate.cost.dominates(&c.cost));
        self.candidates.push(candidate);
        true
    }

    /// Returns the total number of non-dominated candidates on the frontier.
    pub fn len(&self) -> usize {
        self.candidates.len()
    }

    pub fn is_empty(&self) -> bool {
        self.candidates.is_empty()
    }

    /// Selects the optimal candidate according to the specified objective.
    pub fn select_optimal(&self, objective: MrscObjective) -> Option<&OracleCandidate> {
        let cost_model = MrscCostModel::new();
        self.candidates.iter().min_by(|a, b| {
            let score_a = cost_model.score(&a.cost, objective);
            let score_b = cost_model.score(&b.cost, objective);
            score_a.partial_cmp(&score_b).unwrap_or(std::cmp::Ordering::Equal)
        })
    }
}

/// Exhaustive MRSC IDDFS Oracle Engine.
pub struct MrscOracleEngine<'a> {
    func: &'a MirFunction,
    program_funcs: &'a [MirFunction],
    config: IddfsOracleConfig,
    cost_model: MrscCostModel,
}

impl<'a> MrscOracleEngine<'a> {
    pub fn new(
        func: &'a MirFunction,
        program_funcs: &'a [MirFunction],
        config: IddfsOracleConfig,
    ) -> Self {
        Self {
            func,
            program_funcs,
            config,
            cost_model: MrscCostModel::new(),
        }
    }

    /// Run Iterative Deepening Depth-First Search (IDDFS) exploring deepening process trees
    /// from `min_depth` to `max_depth`.
    pub fn explore_iddfs(&self) -> (OracleParetoFrontier, OracleCandidate) {
        let mut frontier = OracleParetoFrontier::new();

        // Include baseline candidate
        let baseline_tree = SupercompilerDriver::new(self.func)
            .with_program_functions(self.program_funcs)
            .run();
        let baseline_cost = self.cost_model.evaluate(self.func, &baseline_tree);
        let baseline_score = self.cost_model.score(&baseline_cost, self.config.objective);
        frontier.insert(OracleCandidate {
            depth: 0,
            residual: self.func.clone(),
            tree: baseline_tree,
            cost: baseline_cost,
            fitness_score: baseline_score,
        });

        // Iterative deepening loop over depths
        let mut current_depth = self.config.min_depth;
        while current_depth <= self.config.max_depth {
            frontier.max_depth_evaluated = frontier.max_depth_evaluated.max(current_depth);
            frontier.total_evaluated += 3;


\end{lstlisting}
